\documentclass[sigconf,anonymous,review]{../acmart/acmart}

\setcopyright{none}

\usepackage{booktabs}
\usepackage{multirow}
\usepackage{amsmath}
\usepackage{array}
\usepackage{microtype}
\usepackage{placeins}

\newcommand{\method}{\textsc{UMER}}
\newcommand{\jsd}{\operatorname{JSD}}
\newcommand{\softmax}{\operatorname{softmax}}
\newcommand{\norm}{\operatorname{norm}}
\newcommand{\mean}{\operatorname{mean}}
\newcommand{\maxpool}{\operatorname{max}}

\acmConference[WSDM '27]
  {the Twentieth ACM International Conference on Web Search and Data Mining}
  {February 15--19, 2027}
  {Hong Kong, China}
\acmBooktitle{Proceedings of the Twentieth ACM International Conference
  on Web Search and Data Mining (WSDM '27), February 15--19, 2027,
  Hong Kong, China}
\acmYear{2027}
\acmDOI{}
\acmISBN{}

\begin{document}

\title[UMER: Unified Multi-Evidence Learning for Rumor Detection]
{UMER: Unified Multi-Evidence Learning for Social Rumor Detection with
Cross-View Consistency and Cross-Event Context}

\begin{abstract}
Social rumor detection remains challenging because seemingly credible source
content, sparse or noisy propagation structures, and incomplete reply
observations may provide conflicting evidence. Existing methods often model
semantic and propagation cues separately or combine independently trained
predictors after training, limiting their ability to learn complementary
evidence jointly. We propose UMER, a unified multi-evidence learning framework
that jointly encodes event propagation and textual semantics while incorporating
cross-event semantic context to complement event-local evidence. Cross-view
consistency encourages two stochastic yet label-preserving serializations of
the same conversation to produce consistent predictive distributions, while
branch-level supervision prevents the unified fusion module from discarding
either structural or textual evidence. To stabilize this heterogeneous
objective under gradient accumulation, we introduce an exact-replay
sharpness-aware optimization procedure that prevents stochastic view changes
from contaminating sharpness estimation. Under fixed five-fold cross-validation
on PHEME, UMER achieves an average accuracy of 0.8988, surpassing the strongest
reproduced baseline by 2.97 percentage points. On Fold 5, where UMER obtains its
highest accuracy of 0.9113, it exceeds the same baseline on the corresponding
test fold by 4.69 percentage points. On Ma-Weibo, UMER reaches 0.9689 average
accuracy, 1.20 percentage points above the strongest completed five-fold
baseline. UMER produces a single jointly optimized checkpoint rather than
relying on post-hoc combinations of independently trained predictors.
\end{abstract}

\begin{CCSXML}
<ccs2012>
 <concept>
  <concept_id>10002951.10003227.10003236</concept_id>
  <concept_desc>Information systems~Social networks</concept_desc>
  <concept_significance>500</concept_significance>
 </concept>
</ccs2012>
\end{CCSXML}

\ccsdesc[500]{Information systems~Social networks}

\keywords{rumor detection, misinformation, graph representation learning,
cross-event context, cross-view consistency, sharpness-aware minimization}

\maketitle

\section{Introduction}

Social media has become an important channel for following breaking events,
sharing eyewitness reports, and organizing public discussion. Its low
publication barrier, however, also allows unverified claims to reach large
audiences before professional verification becomes available. Large-scale
empirical evidence shows that false information can diffuse farther, faster,
deeper, and more broadly than verified truth~\cite{vosoughi2018spread}.
Recent evidence also cautions that exposure is not uniformly distributed:
the most consequential risks are concentrated among highly exposed groups
and in the tail of the consumption distribution~\cite{budak2024harms}.
This combination of rapid diffusion and uneven risk makes timely event-level
detection important, particularly before high-impact cascades become
established.

Much of the evidence needed for detection emerges from the surrounding
conversation. Replies may support, question, repeat, or supplement a claim,
while reply links reveal whether reactions form a broad response around the
source or develop into deeper discussion chains. User and temporal
information further characterize who participates and how the conversation
evolves. The PHEME study shows that users may support an unresolved claim
regardless of its eventual veracity and that informative behavioral signals
arise from the conversational thread as a whole~\cite{zubiaga2016pheme}.
These evidence sources can nevertheless conflict: a linguistically plausible
source may exhibit an unusual propagation pattern, whereas a sparse event
may contain too little structural evidence to support a reliable decision.
Such heterogeneity makes event-level representation particularly difficult
across conversations with substantially different sizes and structures.

Research has progressed from feature-engineered classifiers to sequential
neural models, propagation-aware graph networks, and more recent augmentation
and reconstruction objectives~\cite{castillo2011credibility,ma2016rnn,
ma2018rvnn,bian2020bigcn,cui2024ragcl,wang2025adaat,
wang2026seegraphmae}. Feature-engineered systems depend on predefined
statistics, sequential models can obscure branching reply relations, and
graph models remain bounded by the observed conversation structure.
Augmentation and reconstruction improve robustness but cannot supply
genuinely observed evidence when an event is structurally sparse. These
limitations motivate jointly considering stable event semantics, propagation
structure, and information available across related events.

Pre-trained language models provide strong complementary capacity for event
semantics. DeBERTa disentangles content and positional representations, while
DeBERTaV3 improves pre-training through replaced-token detection and
gradient-disentangled embedding sharing~\cite{he2021deberta,
he2023debertav3}. A social-media event, however, generally contains more text
than can be retained within a fixed input budget. Its representation
therefore depends on which replies are selected, how they are ordered, and
where truncation occurs; two label-preserving serializations of the same
conversation can consequently produce different predictions. In addition,
encoding each event independently does not explicitly reveal whether it is
semantically close to previously observed rumor or non-rumor patterns.
Sentence-level representations make such semantic neighborhoods
measurable~\cite{reimers2019sbert}, suggesting that evidence from related
events can complement event-local text and propagation signals.

Bringing structural evidence, multiple text views, and cross-event semantic
support into one detector introduces a further optimization challenge. These
components are learned through multiple supervised objectives and stochastic
event representations, whose variation becomes especially important under
gradient accumulation. Prediction-level consistency methods such as
R-Drop~\cite{liang2021rdrop} show that stochastic views can be regularized by
encouraging agreement between their oant
node-content encoding with a structural Transformer to represent
conversations of varying sizes and topologies. The event-text channel encodes
two label-preserving serializations using a shared DeBERTa-v3-large encoder
and promotes agreement between their predictive distributions. The
cross-event semantic channel estimates local class support from related
reference events, providing additional evidence for ambiguous or structurally
sparse conversations. A unified fusion head integrates the three decision
signals, while branch-level supervision maintains the discriminative capacity
of the parametric channels. Finally, exact-replay SAM preserves the same
accumulated events, constructed views, and stochastic states across its two
optimization passes.

Our main contributions are threefold:

\begin{itemize}
  \item We propose a unified multi-evidence rumor detection architecture that
  jointly models propagation structure, event-level language, and cross-event
  semantic context within a single trainable model. The cross-event channel
  provides explicit semantic support for events whose local textual or
  structural evidence is insufficient.

  \item We introduce a robust evidence-learning mechanism that combines
  complementary graph readouts, label-preserving event serializations,
  prediction-level cross-view consistency, and branch supervision. This
  design improves stability under reply selection, ordering, and truncation
  while maintaining the individual competence of the structural and textual
  channels.

  \item We develop an exact effective-batch SAM procedure for accumulated
  multi-view training. By replaying event identities, constructed views, and
  stochastic states, the procedure preserves the intended
  parameter-neighborhood objective across the two SAM passes.
\end{itemize}

\section{Related Work}

Research related to our work spans three connected questions: how to represent
heterogeneous evidence within an event, how to obtain useful context beyond
the event itself, and how to optimize the resulting stochastic multi-objective
model reliably. We therefore review event content and propagation modeling,
cross-event semantic context, and consistency and sharpness-aware learning in
sequence.

\subsection{Event Content, Propagation, and Joint Modeling}

Early rumor detectors used manually constructed linguistic, user, temporal,
and diffusion features~\cite{castillo2011credibility,ma2015svmts}. Recurrent
models subsequently learned temporal representations directly from sequences
of posts~\cite{ma2016rnn}, reducing dependence on predefined statistics but
flattening the branching relations among replies. Pre-trained language models
provide substantially stronger semantic representations. In particular,
DeBERTa separates content and positional representations, and DeBERTaV3
improves its pre-training objective through replaced-token detection and
gradient-disentangled embedding sharing~\cite{he2021deberta,
he2023debertav3}. Nevertheless, a language encoder still receives a finite
serialization of an event, making its output sensitive to which replies are
retained and how a branching conversation is ordered.

Propagation models preserve information that is lost during such
linearization. Bi-GCN applies graph convolution in top-down and bottom-up
directions to capture both claim propagation and response
dispersion~\cite{bian2020bigcn}. RAGCL adapts graph contrastive augmentation
to the wide and shallow structures observed in rumor
trees~\cite{cui2024ragcl}. Temporal propagation optimization and FGDGNN
further model time intervals and graph evolution within and across
snapshots~\cite{peng2025temporal,guo2025fgdgnn}. Robustness-oriented methods
construct adaptive or adversarial graph views~\cite{wang2025adaat}, or recover
masked local and global semantic-evolution patterns~\cite{wang2026seegraphmae}.
These advances improve propagation representation, but the available signal
remains bounded by the observed event graph and is consequently weak for
source-only or extremely sparse conversations.

Several methods jointly learn content and structural evidence. Rumor2vec
combines source-text encoding with representations learned from a union graph
that links the propagation structures of multiple events through shared
network information~\cite{tu2021rumor2vec}. DUCK uses Transformers and graph
attention networks to jointly model conversation content, comment
propagation, and participating-user networks~\cite{tian2022duck}. TGA-CMF
combines time-decayed graph attention with graph-guided textual fusion for
early rumor detection~\cite{ding2026tgacmf}. These methods establish that
end-to-end content--structure learning is effective and that joint training
itself is not a new contribution. Their evidence, however, is primarily
derived from the observed conversation, user connectivity, or temporal graph;
they do not simultaneously address instability across valid event-text
serializations and instance-level semantic support from related training
events.

\subsection{Cross-Event Semantic Context}

Event-local encoders cannot directly express how a new conversation relates
to previously observed events. Sentence-BERT demonstrates that semantically
related texts can be organized into efficient cosine
neighborhoods~\cite{reimers2019sbert}, but semantic similarity alone does not
specify how information should be transferred between rumor events.

Existing rumor studies use cross-event information mainly for transfer and
generalization. Cross-topic rumor detection represents training topics and
combines them according to the topic distribution of a target
instance~\cite{ren2021crosstopic}. Few-shot cross-event detection instead
learns class prototypes or relation metrics from historical events through
episodic training and adapts them to novel events with a small number of
labeled examples~\cite{ran2023metacross}. In an adjacent misinformation
setting, cross-document event graph reasoning aggregates complementary and
contradictory information from clusters of related news
documents~\cite{wu2022crossdoc}. Rumor2vec also shares information across
events, but does so through user and propagation connectivity rather than
instance-level semantic neighborhoods~\cite{tu2021rumor2vec}.

These approaches demonstrate the value of knowledge beyond an isolated event,
but they target cross-topic transfer, few-shot adaptation, clustered external
documents, or shared propagation networks. They leave open how to expose
instance-level semantic precedent under ordinary event-level supervised
evaluation while keeping the admissible reference set explicit and isolated
from validation and test labels.

\subsection{Consistency and Sharpness-Aware Optimization}

Consistency regularization stabilizes predictions under stochastic
perturbations. R-Drop evaluates the same input twice with independently
sampled dropout masks and minimizes supervised loss together with a
bidirectional KL divergence between the two output
distributions~\cite{liang2021rdrop}. It therefore regularizes model-internal
dropout variation, but does not directly address input-level changes caused by
reply selection, ordering, or truncation. For two complete event
serializations, Jensen--Shannon divergence offers a symmetric and bounded
distribution-level constraint through their mean predictive distribution.
The corresponding event-view construction and objective are specified in
Section~\ref{sec:text}.

Sharpness-Aware Minimization searches for parameters that retain low loss
under a bounded parameter perturbation~\cite{foret2021sam}. Subsequent
variants refine different aspects of this neighborhood. ASAM makes sharpness
less sensitive to parameter rescaling by adapting the perturbation region to
parameter magnitude~\cite{kwon2021asam}. GSAM explicitly reduces a surrogate
gap associated with local curvature in addition to the perturbed
loss~\cite{zhuang2022gsam}, whereas Fisher SAM replaces the Euclidean ball
with an ellipsoid induced by Fisher information
geometry~\cite{kim2022fishersam}. These variants improve the geometry or
objective used to define sharpness, but they retain the requirement that the
unperturbed and perturbed passes evaluate the same stochastic training
objective.

We use the original Euclidean SAM objective as a controlled base because it
introduces neither layer-wise scaling, Fisher estimation, nor an additional
surrogate-gradient decomposition. This choice isolates the issue addressed in
our optimization design: preserving the exact effective batch and stochastic
event realization across the two passes. Exact replay is orthogonal to the
choice of neighborhood geometry and could in principle also be applied to
adaptive or information-geometric SAM variants.

Taken together, prior work provides strong event-local semantic and
propagation encoders, mechanisms for cross-topic or cross-document transfer,
and generic consistency or sharpness-aware regularization. Our work targets
their intersection: one jointly optimized detector combines event-internal
propagation and language evidence with explicit training-event semantic
support, regularizes complete event serializations at the prediction level,
and preserves the intended sharpness objective under accumulated stochastic
training.

\section{Methodology}
\label{sec:method}

This section presents \method\ from task definition to joint optimization.
Section~\ref{sec:problem} formalizes event-level rumor detection and summarizes
the notation.  Sections~\ref{sec:graph}--\ref{sec:context} describe the three
evidence channels: propagation, cross-view text, and cross-event context.
Section~\ref{sec:fusion} introduces decision-level fusion and its supervised
objective, and Section~\ref{sec:sam} gives the exact-replay SAM procedure.

\subsection{Problem Formulation}
\label{sec:problem}

Let an event be
\begin{equation}
  e_i=(G_i,X_i,S_i,y_i),
  \qquad G_i=(V_i,E_i), \quad y_i\in\{0,1\},
  \label{eq:event}
\end{equation}
where $G_i$ is a directed conversation graph, $V_i$ contains the source post
and its replies, $E_i$ preserves reply relations, $X_i$ denotes node-level
semantic and social attributes, and $S_i$ is a source-and-replies text
serialization.  We use $y_i=1$ for rumor and $y_i=0$ for non-rumor.  Given an
event set $\mathcal{D}$ and an admissible reference event bank $\mathcal{B}$,
the objective is to learn one classifier
\begin{equation}
  f_{\theta}(e_i,\mathcal{B})
  =\boldsymbol{p}_i\in[0,1]^2,\qquad
  \sum_{c\in\{0,1\}}p_{i,c}=1.
  \label{eq:classifier}
\end{equation}
Unlike a post-hoc ensemble, all parametric encoders and the fusion head are
optimized as one model.  The bank supplies a deterministic context feature
and does not introduce an independently selected neural checkpoint.

\begin{table}[H]
  \caption{Main notation used in the methodology.}
  \label{tab:notation}
  \centering
  \small
  \begin{tabular}{@{}p{0.20\columnwidth}p{0.71\columnwidth}@{}}
    \toprule
    Symbol & Meaning \\
    \midrule
    $e_i$ & Event indexed by $i$ \\
    $G_i=(V_i,E_i)$ & Directed conversation graph \\
    $X_i,S_i$ & Node attributes and serialized event text \\
    $y_i$ & Binary rumor label \\
    $\boldsymbol{h}_i^g$ & Propagation representation \\
    $\boldsymbol{z}_i^g$ & Two-class graph evidence logits \\
    $S_i^{(a)}$ & Stochastic text view $a$ \\
    $\boldsymbol{z}_i^{t,a}$ & Text logits for view $a$ \\
    $\bar{\boldsymbol{z}}_i^t$ & Mean two-view text evidence \\
    $\boldsymbol{q}_i$ & Normalized event descriptor \\
    $\mathcal{B}$ & Reference event bank \\
    $\mathcal{N}_k(i)$ & Legal $k$-neighbor set for event $i$ \\
    $r_i$ & Cross-event semantic context margin \\
    $\boldsymbol{z}_i$ & Final fused logits \\
    $\theta,\rho$ & Model parameters and SAM radius \\
    \bottomrule
  \end{tabular}
\end{table}

The conversation may contain only a source or a small number of replies.
Accordingly, every channel is defined for a nonempty source event: the graph
readout uses the source when no reply is present, the text serializer always
retains the source, and the event descriptor has a valid source, mean, and
maximum.  This design prevents sparse conversations from being discarded
during preprocessing and lets cross-event context supplement weak within-event
evidence.

\begin{figure*}[t]
  \centering
  \includegraphics[width=\textwidth]{figures/fig1_umer_architecture.pdf}
  \caption{Overview of UMER.  An event is represented by three complementary
  modules: a two-channel propagation encoder, cross-view textual consistency
  with a shared DeBERTa-v3 encoder, and cross-event semantic context.  Their
  graph, text, and context evidence is jointly fused to produce the rumor
  prediction.}
  \Description{A black-and-white left-to-right architecture diagram.  A reply
  topology and its textual content feed three complementary modules: a
  two-channel propagation encoder, two label-preserving text views processed
  by a shared encoder, and query-to-memory semantic retrieval.  The resulting
  graph evidence, text evidence, and context support enter a joint evidence
  fusion module that outputs the rumor prediction.}
  \label{fig:architecture}
\end{figure*}

\subsection{Propagation Representation Learning}
\label{sec:graph}

Propagation structure provides evidence that is not reducible to wording.
Rumor and non-rumor discussions may differ in breadth, response timing, and
the interaction between the source and replies.  A single graph readout can
nevertheless be brittle.  Attention over node sequences is expressive but may
overfit the ordering of a small node set, whereas global pooling is stable but
can erase important interactions.  We therefore combine a
permutation-invariant token-set path with a structural Transformer path.

\paragraph{Node-token set path.}
For node $v\in V_i$, let
$\boldsymbol{u}_{ivm}\in\mathbb{R}^{d_m}$ be its $m$th non-text token,
covering user, temporal, and propagation summaries.  A shared multilayer
perceptron $\phi$ projects the $M=14$ tokens:
\begin{equation}
  \boldsymbol{o}_{ivm}
  =\phi(\boldsymbol{u}_{ivm}),\qquad
  \boldsymbol{s}_{iv}
  =[\mean_m\boldsymbol{o}_{ivm};
    \maxpool_m\boldsymbol{o}_{ivm}].
  \label{eq:token-set}
\end{equation}
Mean pooling captures the overall evidence level and max pooling preserves a
salient structural cue.  A second readout over nodes yields
\begin{equation}
  \boldsymbol{h}_i^{\mathrm{set}}
  =[\mean_{v\in V_i}\boldsymbol{s}_{iv};
    \maxpool_{v\in V_i}\boldsymbol{s}_{iv}],
  \label{eq:set-readout}
\end{equation}
which is invariant to node enumeration and remains valid for a one-node
event.

\paragraph{Structural Transformer path.}
Let $\boldsymbol{x}_{iv}\in\mathbb{R}^{384}$ be a fixed sentence embedding and
$\boldsymbol{a}_{iv}$ collect non-text and temporally discretized structural
features.  Separate projections place them in a common $d$-dimensional space:
\begin{equation}
  \boldsymbol{n}_{iv}
  =W_x\boldsymbol{x}_{iv}+W_a\boldsymbol{a}_{iv}+\boldsymbol{b}_n.
  \label{eq:node-projection}
\end{equation}
A learned classification token is prepended and a one-layer, eight-head
Transformer encodes the variable node set:
\begin{align}
  H_i &=\operatorname{Transformer}
  ([\boldsymbol{c};\boldsymbol{n}_{i1};\ldots;
  \boldsymbol{n}_{i|V_i|}]),\\
  \boldsymbol{h}_i^{\mathrm{tr}}
  &=[H_{i,0};\mean_{v\in V_i}H_{i,v}].
  \label{eq:graph-transformer}
\end{align}
The classification token aggregates learned interactions, while the explicit
mean creates a direct global path that is less dependent on a single token.
The two graph summaries are fused and classified:
\begin{equation}
  \boldsymbol{h}_i^g
  =F_g([\boldsymbol{h}_i^{\mathrm{set}};
        \boldsymbol{h}_i^{\mathrm{tr}}]),\qquad
  \boldsymbol{z}_i^g=W_g\boldsymbol{h}_i^g+\boldsymbol{b}_g.
  \label{eq:graph-logits}
\end{equation}
Mild source-preserving node dropout perturbs the reply context during
optimization.  The source node is never removed, so the branch remains
well-defined for sparse threads.

\subsection{Cross-View Textual Representation Learning}
\label{sec:text}

An event contains far more text than a conventional sentence classifier can
always retain.  A deterministic serialization may overemphasize whichever
replies happen to appear before the token limit.  We instead sample two legal
views $S_i^{(1)}$ and $S_i^{(2)}$ from a view distribution
$\mathcal{T}(S_i)$.  Both views preserve the source and use only posts from
the same event, but may expose different reply orderings and truncation
boundaries:
\begin{equation}
  S_i^{(a)}\sim\mathcal{T}(S_i),\qquad a\in\{1,2\}.
  \label{eq:text-views}
\end{equation}
This construction treats reply selection as an internal observation
uncertainty rather than creating new labeled examples.

\paragraph{Shared semantic encoder.}
Both views pass through the same DeBERTa-v3-large encoder $T_{\theta_t}$ and
classification head:
\begin{equation}
  \boldsymbol{z}_i^{t,a}
  =W_tT_{\theta_t}(S_i^{(a)})+\boldsymbol{b}_t,\qquad
  \boldsymbol{p}_i^{t,a}=\softmax(\boldsymbol{z}_i^{t,a}).
  \label{eq:text-logits}
\end{equation}
Sharing weights ensures that the views regularize one text model rather than
form two independently selected predictors.  Their decision evidence is
\begin{equation}
  \bar{\boldsymbol{z}}_i^t
  =\tfrac{1}{2}
  (\boldsymbol{z}_i^{t,1}+\boldsymbol{z}_i^{t,2}).
  \label{eq:text-average}
\end{equation}

\paragraph{Prediction consistency.}
Let
$\boldsymbol{m}_i=\tfrac{1}{2}
(\boldsymbol{p}_i^{t,1}+\boldsymbol{p}_i^{t,2})$.
We regularize the two decisions using the Jensen--Shannon
divergence
\begin{equation}
  \mathcal{L}_{\mathrm{cv}}^{(i)}
  =\tfrac{1}{2}\operatorname{KL}
  (\boldsymbol{p}_i^{t,1}\Vert\boldsymbol{m}_i)
  +\tfrac{1}{2}\operatorname{KL}
  (\boldsymbol{p}_i^{t,2}\Vert\boldsymbol{m}_i).
  \label{eq:gjsd}
\end{equation}
The constraint acts on class distributions, not hidden vectors.  Thus the
views can encode different reply details while being discouraged from
contradicting each other at the decision level.

\subsection{Cross-Event Semantic Context Learning}
\label{sec:context}

Within-event evidence can be incomplete.  A source-only event has little
propagation information, and a short generic claim may remain ambiguous even
to a large language model.  Cross-event context asks a complementary
question: do semantically related reference events support one class more
strongly than the other?  The resulting score is not a replacement for graph
or text encoding; it is an explicit local precedent whose neighbors and
weights can be audited.

\paragraph{Event descriptor.}
Let $\boldsymbol{x}_{iv}$ be the fixed 384-dimensional sentence embedding of
post $v$, and let $s_i$ identify the source.  We summarize complementary
aspects of the event using
\begin{equation}
 \boldsymbol{q}_i=\norm\!\left[
 \boldsymbol{x}_{is_i};
 \mean_{v}\boldsymbol{x}_{iv};
 \maxpool_{v}\boldsymbol{x}_{iv};
 \mean_{v}\boldsymbol{x}_{iv}-\boldsymbol{x}_{is_i}
 \right].
 \label{eq:event-descriptor}
\end{equation}
The source represents the original claim, the mean captures overall
discussion semantics, the maximum preserves dimension-wise salient signals,
and the source-relative term describes how replies shift the conversation.
For a source-only event, the relative term becomes zero and the remaining
components remain valid.

\paragraph{Reference bank and local support.}
Define the reference event bank as
\begin{equation}
  \mathcal{B}=\{(\boldsymbol{q}_j,y_j)\}_{j\in\mathcal{I}_{\mathcal{B}}}.
  \label{eq:bank}
\end{equation}
For query $i$, its legal $k$-nearest neighbor set
$\mathcal{N}_k(i)$ excludes the query itself whenever it is indexed in the
bank.  With cosine similarity and temperature $\tau$,
\begin{align}
  a_{ij}
  &=\frac{\exp(\cos(\boldsymbol{q}_i,\boldsymbol{q}_j)/\tau)}
  {\sum_{\ell\in\mathcal{N}_k(i)}
   \exp(\cos(\boldsymbol{q}_i,\boldsymbol{q}_{\ell})/\tau)},\\
  p_i^m&=\sum_{j\in\mathcal{N}_k(i)}a_{ij}y_j.
  \label{eq:neighbor-support}
\end{align}
The semantic context is expressed as a log-odds margin
\begin{equation}
  r_i=\log
  \frac{p_i^m+\epsilon}{1-p_i^m+\epsilon}.
  \label{eq:context-margin}
\end{equation}
A positive margin supports the rumor class and a negative margin supports the
non-rumor class.  Temperature-weighted support is smoother than copying the
single nearest label, while the scalar margin prevents the non-parametric
channel from dominating by dimensionality alone.  The admissible membership
of $\mathcal{B}$ is specified by the evaluation protocol in
Section~\ref{sec:data}.

\subsection{Decision-Level Tri-Evidence Fusion}
\label{sec:fusion}

High-dimensional hidden-state concatenation can introduce a large fusion
network on top of a modest event corpus.  It can also obscure whether a
prediction arose from propagation, text, or cross-event support.  We instead
fuse decision-level evidence: two graph logits, two averaged text
logits, and one semantic context margin,
\begin{equation}
  \boldsymbol{u}_i=
  [\boldsymbol{z}_i^g;\bar{\boldsymbol{z}}_i^t;r_i]
  \in\mathbb{R}^{5}.
  \label{eq:fusion-input}
\end{equation}
The final classifier is deliberately linear:
\begin{equation}
  \boldsymbol{z}_i=W_f\boldsymbol{u}_i+\boldsymbol{b}_f,\qquad
  \hat{\boldsymbol{p}}_i=\softmax(\boldsymbol{z}_i).
  \label{eq:fusion}
\end{equation}
Although compact, $W_f$ can learn class-specific agreement and conflict
patterns across channels.  The entire path from graph and text encoders to
$W_f$ is optimized jointly, producing one deployable checkpoint.

Class imbalance is handled by weighted cross entropy
$\operatorname{CE}_w$.  Auxiliary graph and text losses preserve the
discriminative ability of each parametric branch:
\begin{align}
  \mathcal{L}_{\mathrm{sup}}^{(i)}
  ={}&\operatorname{CE}_w(\boldsymbol{z}_i,y_i)
  +\lambda_g\operatorname{CE}_w(\boldsymbol{z}_i^g,y_i)\nonumber\\
  &+\lambda_t
  \operatorname{CE}_w(\bar{\boldsymbol{z}}_i^t,y_i).
  \label{eq:supervised-loss}
\end{align}
The complete per-event objective is
\begin{equation}
  \mathcal{L}^{(i)}
  =\mathcal{L}_{\mathrm{sup}}^{(i)}
  +\lambda_{\mathrm{cv}}\mathcal{L}_{\mathrm{cv}}^{(i)}.
  \label{eq:complete-loss}
\end{equation}
Branch supervision is important because a low-capacity fusion head can
otherwise settle on the easiest channel early and starve the complementary
encoder of a useful learning signal.

\subsection{Exact Effective-Batch Sharpness-Aware Optimization}
\label{sec:sam}

\paragraph{Why exact replay is necessary.}
SAM minimizes the maximum loss within a radius-$\rho$ parameter neighborhood:
\begin{equation}
  \min_{\theta}
  \max_{\|\boldsymbol{\epsilon}\|_2\le\rho}
  \mathcal{L}_{B}(\theta+\boldsymbol{\epsilon}).
  \label{eq:sam-objective}
\end{equation}
Using the first-order approximation,
\begin{equation}
  \boldsymbol{\epsilon}_{B}
  =\rho\,
  \frac{\nabla_{\theta}\mathcal{L}_{B}(\theta)}
  {\|\nabla_{\theta}\mathcal{L}_{B}(\theta)\|_2+\delta}.
  \label{eq:sam-perturbation}
\end{equation}
The notation $B$ must refer to the same stochastic objective in both passes.
If the perturbed pass samples different event views, node dropout, or encoder
dropout, then
$\mathcal{L}_{B'}(\theta+\boldsymbol{\epsilon}_{B})$ is evaluated for
$B'\ne B$.  The update then mixes parameter sharpness with view noise.

\paragraph{Accumulated two-pass update.}
Let the effective batch be the ordered union
$B=\bigcup_{m=1}^{A}B_m$ of $A$ physical microbatches.  For each optimizer
step, \method\ executes:
\begin{align}
  \boldsymbol{g}
  &=\frac{1}{A}\sum_{m=1}^{A}
    \nabla_{\theta}\mathcal{L}_{B_m}(\theta;\omega_m),
  \label{eq:first-pass}\\
  \boldsymbol{g}_{\mathrm{sam}}
  &=\frac{1}{A}\sum_{m=1}^{A}
    \nabla_{\theta}
    \mathcal{L}_{B_m}(\theta+\boldsymbol{\epsilon}_{B};\omega_m),
  \label{eq:second-pass}
\end{align}
where $\omega_m$ contains the constructed graph/text views and random-number
generator state for microbatch $m$.  The same $\omega_m$ appears in both
equations.  In practice, event identities, views, and random states are cached
during the first pass; parameters are perturbed once after all first-pass
gradients have accumulated; the microbatches are replayed in the same order;
parameters are restored; and AdamW applies only the clipped perturbed-pass
gradient.  SAM adds no parameters and no inference-time operations.

\section{Experiments}
\label{sec:experiments}

\subsection{Datasets}
\label{sec:data}

We evaluate on PHEME and Ma-Weibo, two public event-level rumor datasets from
different platforms and languages.  Table~\ref{tab:datasets} reports the
class composition and the post-count distribution after our common
preprocessing, i.e., the instances actually supplied to the models.  PHEME
contains journalist-verified English Twitter conversation threads around nine
breaking-news events~\cite{zubiaga2016pheme}.  Its processed binary corpus has
6,425 event graphs, of which 2,402 are rumors and 4,023 are non-rumors.
Ma-Weibo contains Chinese Sina Weibo events collected from rumors reported by
the platform's community management center and a similar number of unreported
general events~\cite{ma2016rnn}.  It has 4,664 events, including 2,313 rumors
and 2,351 non-rumors.  The original release reports 3,805,656 posts, whereas
the common observation window below retains 1,771,329 posts.

\begin{table}[t]
  \caption{Dataset statistics after the common 240-hour preprocessing.
  ``Posts/event'' counts the retained source post and replies or reposts in an
  event; Q1 and Q3 are the first and third quartiles.}
  \label{tab:datasets}
  \centering
  \small
  \setlength{\tabcolsep}{6pt}
  \begin{tabular}{@{}lrr@{}}
    \toprule
    Statistic & PHEME & Ma-Weibo \\
    \midrule
    Language & English & Chinese \\
    Events & 6,425 & 4,664 \\
    Rumor / non-rumor & 2,402 / 4,023 & 2,313 / 2,351 \\
    Retained posts & 105,238 & 1,771,329 \\
    Posts/event minimum & 1 & 1 \\
    Posts/event mean $\pm$ SD & 16.38 $\pm$ 19.59 &
    379.79 $\pm$ 340.78 \\
    Posts/event median [Q1, Q3] & 12 [5, 21] & 236 [112, 592] \\
    Posts/event maximum & 342 & 1,021 \\
    Source-only events & 677 (10.54\%) & 6 (0.13\%) \\
    \bottomrule
  \end{tabular}
\end{table}

Both corpora are converted to a common event representation.  The source post
is the root, and retained replies or reposts become event-local nodes ordered
chronologically.  Parent--child relations are preserved when available;
records with missing parent identifiers are attached using the dataset's
recoverable propagation metadata rather than discarded.  Text is normalized
without introducing external articles, user profiles, generated labels, or
additional rumor corpora.  Each node receives its post-semantic vector,
available user attributes, relative time, and propagation summaries.  Posts
within 240 hours of the source are retained and time is discretized into 480
half-hour bins, providing the same observation window on both datasets.

The two corpora differ sharply in conversation density.  A processed PHEME
event contains only 16.38 posts on average, compared with 379.79 in Ma-Weibo,
and their strongly right-skewed distributions make the medians and quartiles
more informative than the maxima alone.  In particular, 677 PHEME events
(10.54\%) contain only the source post, whereas this occurs for just six
Ma-Weibo events (0.13\%).  We retain these sparse cases: the source is always
preserved under text truncation and node dropout, and an empty reply aggregate
falls back to the source-defined representation in
Section~\ref{sec:method}.  Removing them would change the benchmark population
and conceal precisely the cases in which within-event propagation is least
informative and cross-event context may be most useful.

For comparability, both datasets use fixed stratified five-fold event-level
partitions.  All posts belonging to one event remain in the same fold.  Within
each outer training portion, a validation subset is separated before
optimization.  The reference event bank is populated only with labeled events
available to the corresponding outer training partition; validation and test
events query that bank but never populate it, and an indexed training query is
excluded from its own neighborhood.  Thus graph construction, text views, and
cross-event context respect the same event boundary and cannot transfer labels
from the held-out fold.

\subsection{Baselines and Evaluation Metrics}
\label{sec:baselines}

\paragraph{Baselines.}
We compare against nine methods spanning feature engineering, sequential
modeling, graph neural networks, and recent robust or temporal architectures.
\textbf{SVM-TS} represents an event by time series of content, user, and
propagation statistics and applies a linear SVM~\cite{ma2015svmts}.
\textbf{DLSTM} groups a source and temporally ordered replies, filters reply
content using user influence, and encodes the sequence with stacked
LSTMs~\cite{xu2019dlstm}.  \textbf{Bi-GCN} combines top-down and bottom-up
graph convolution with root enhancement~\cite{bian2020bigcn}.
\textbf{RAGCL} adds centrality-guided node, edge, and attribute augmentations
to graph contrastive learning~\cite{cui2024ragcl}.  \textbf{PAGCN} couples
semantic encoding with propagation-path graph convolution and path
attention~\cite{liu2025pagcn}.  \textbf{TGA-CMF} uses time-decayed graph
attention and graph-guided textual fusion~\cite{ding2026tgacmf}.
\textbf{ADAAT-FGM} combines adaptive edge/feature augmentation, adversarial
hard examples, and supervised contrastive learning~\cite{wang2025adaat}.
\textbf{RRD-N} uses information-bottleneck contrastive pre-training and
Bayesian variational fine-tuning to resist propagation noise
~\cite{zhang2026rrdn}.  Finally, \textbf{MST} constructs cumulative temporal
graphs and fuses dual dynamic graph convolutions through cross-attention
~\cite{huang2025mst}.

Every method receives the same fixed outer five-fold event partitions.
Beyond this common partition, we retain each paper or author implementation's
preprocessing, model, optimizer, and stopping protocol whenever disclosed;
missing settings are reported with the reproduction rather than silently
attributed to the authors.  The released MST pipeline cannot construct a
propagation graph for source-only events and therefore evaluates only its
eligible subset; all other reported systems retain the complete test fold.

\begin{table}[b]
  \caption{Training and evaluation configuration of \method.}
  \label{tab:training-config}
  \centering
  \small
  \setlength{\tabcolsep}{4pt}
  \begin{tabular}{@{}ll@{}}
    \toprule
    Item & Setting \\
    \midrule
    Text encoder & DeBERTa-v3-large \\
    Graph hidden size / dropout & 256 / 0.35 \\
    Physical / effective batch & 4 / 64 \\
    Gradient accumulation & 16 microbatches \\
    Graph / text / fusion LR & $10^{-4}$ / $2{\times}10^{-5}$ / $10^{-3}$ \\
    Weight decay / warmup & 0.01 / 0.10 \\
    Epoch cap / patience & 30 / 7 \\
    Retrieval $k$ / temperature & 31 / 0.05 \\
    Cross-view weight $\lambda_{\mathrm{cv}}$ & 0.5 \\
    SAM radius $\rho$ & 0.05 \\
    Partition / training seed & 3090 / 2000 \\
    Class weighting & Square-root inverse frequency \\
    Checkpoint criterion & Max--min of four validation metrics \\
    \bottomrule
  \end{tabular}
\end{table}

\begin{table*}[!t]
  \caption{Five-fold mean comparison on PHEME.  P, R, and F1 denote
  class-wise precision, recall, and F1, respectively.  The best value in each
  column is shown in bold.}
  \label{tab:pheme-results}
  \centering
  \small
  \setlength{\tabcolsep}{7.2pt}
  \begin{tabular}{@{}lccccccc@{}}
    \toprule
    & \multicolumn{3}{c}{Rumor} &
      \multicolumn{3}{c}{Non-rumor} & Overall \\
    \cmidrule(lr){2-4}\cmidrule(lr){5-7}
    Method & P & R & F1 & P & R & F1 & Acc. \\
    \midrule
    SVM-TS       & 0.4829 & 0.4430 & 0.4620 & 0.6830 & 0.7166 & 0.6994 & 0.6143 \\
    DLSTM        & 0.4147 & 0.4218 & 0.3870 & 0.6636 & 0.6557 & 0.6450 & 0.5682 \\
    Bi-GCN
                  & 0.8390 & 0.7589 & 0.7961 & 0.8642 & 0.9122 & 0.8873 & 0.8549 \\
    RAGCL
                  & 0.7291 & 0.7276 & 0.7272 & 0.8380 & 0.8371 & 0.8370 & 0.7961 \\
    PAGCN         & 0.7734 & 0.6173 & 0.6528 & 0.8053 & 0.8832 & 0.8383 & 0.7838 \\
    TGA-CMF       & 0.6840 & 0.6653 & 0.6727 & 0.8039 & 0.8146 & 0.8085 & 0.7588 \\
    ADAAT-FGM     & 0.7742 & 0.7090 & 0.7371 & 0.8357 & 0.8730 & 0.8528 & 0.8117 \\
    RRD-N
                  & 0.8357 & 0.7919 & 0.8119 & 0.8803 & 0.9061 & 0.8926 & 0.8633 \\
    MST$^{\ddagger}$
                  & \textbf{0.8620} & 0.7760 & 0.8161 & 0.8738 &
                    \textbf{0.9247} & 0.8983 & 0.8691 \\
    \midrule
    \method       & 0.8560 & \textbf{0.8784} & \textbf{0.8666} &
                    \textbf{0.9264} & 0.9110 & \textbf{0.9185} &
                    \textbf{0.8988} \\
    \bottomrule
  \end{tabular}
\end{table*}

\paragraph{Evaluation metrics.}
For each class $c\in\{\mathrm{R},\mathrm{NR}\}$, we treat $c$ as the positive
class and denote its true positives, false positives, and false negatives by
$\mathrm{TP}_c$, $\mathrm{FP}_c$, and $\mathrm{FN}_c$.
\begin{align}
  P_c&=\frac{\mathrm{TP}_c}
  {\mathrm{TP}_c+\mathrm{FP}_c},\qquad
  R_c=\frac{\mathrm{TP}_c}
  {\mathrm{TP}_c+\mathrm{FN}_c},\\
  F_{1,c}&=\frac{2P_cR_c}{P_c+R_c}.
\end{align}
Overall accuracy is
\begin{equation}
  \mathrm{Accuracy}
  =\frac{\mathrm{TP}_{\mathrm{R}}+\mathrm{TP}_{\mathrm{NR}}}
  {n_{\mathrm{R}}+n_{\mathrm{NR}}},
\end{equation}
where $n_c$ is the support of class $c$.  Accuracy is reported only once:
under binary one-versus-rest evaluation, a class-specific accuracy equals
overall accuracy, while correctness restricted to examples of class $c$
equals $R_c$.  Repeating either quantity as a separate class-accuracy column
would therefore duplicate an existing metric.  For every reported quantity,
we first compute the metric on each fixed test fold and then take the
arithmetic mean over the five folds.  We summarize the two labels using
\begin{equation}
  P_{\mathrm{macro}}=\frac{1}{2}\sum_c P_c,\quad
  R_{\mathrm{macro}}=\frac{1}{2}\sum_c R_c,\quad
  F_{1,\mathrm{macro}}=\frac{1}{2}\sum_c F_{1,c}.
\end{equation}
The PHEME table gives the full class-wise decomposition.  The Ma-Weibo table
reports accuracy once and macro-averages precision, recall, and F1 across the
rumor and non-rumor labels, preventing the larger class from receiving more
weight.

\subsection{Experimental Results and Analyses}
\label{sec:results}

\paragraph{PHEME comparison.}
Table~\ref{tab:pheme-results} reports five-fold mean performance on PHEME.
Boldface marks the highest value in each metric column.  The double dagger
identifies the MST result computed after its released pipeline excludes
source-only events.  This qualification is necessary because the corresponding
row is not evaluated over every test event.

\method\ achieves the highest overall accuracy, 0.8988, improving over the
strongest reproduced baseline, MST, by 2.97 percentage points.  The gain is
not obtained by favoring only the majority non-rumor class.  Compared with
MST, \method\ raises rumor recall from 0.7760 to 0.8784 and rumor F1 from
0.8161 to 0.8666, gains of 10.24 and 5.05 percentage points, respectively.
It also improves non-rumor precision from 0.8738 to 0.9264 and non-rumor F1
from 0.8983 to 0.9185.  MST retains slightly higher rumor precision and
non-rumor recall, indicating a more conservative tendency to assign the rumor
label.  In contrast, the unified evidence learned by \method\ substantially
reduces missed rumors while preserving strong performance on non-rumors.

Among the remaining baselines, RRD-N and Bi-GCN obtain overall accuracies of
0.8633 and 0.8549.  The sequence-based DLSTM and feature-engineered SVM-TS are
markedly weaker, whereas graph-oriented methods generally perform better,
supporting the value of propagation structure.  Nevertheless, their lower
rumor recall and F1 show that event-local propagation alone is insufficient
for the minority class.  The class-wise gains of \method\ are consistent with
its joint use of propagation, cross-view textual evidence, and cross-event
semantic support.

\begin{table}[t]
  \caption{Results on Ma-Weibo. P, R, and F1 are two-class macro-averages.
  All rows are five-fold means except RRD-N$^{\dagger}$, which is the mean of
  the three completed folds. The incomplete row is excluded from five-fold
  ranking claims.}
  \label{tab:maweibo-results}
  \centering
  \small
  \setlength{\tabcolsep}{3.6pt}
  \begin{tabular}{@{}lcccc@{}}
    \toprule
    Method & Acc. & P & R & F1 \\
    \midrule
    SVM-TS       & 0.7706 & 0.7714 & 0.7704 & 0.7703 \\
    DLSTM        & 0.5967 & 0.5984 & 0.5969 & 0.5950 \\
    Bi-GCN       & 0.8986 & 0.8989 & 0.8985 & 0.8985 \\
    RAGCL        & 0.9117 & 0.9195 & 0.9116 & 0.9109 \\
    PAGCN        & 0.9569 & 0.9570 & 0.9570 & 0.9569 \\
    TGA-CMF      & 0.9031 & 0.9062 & 0.9034 & 0.9029 \\
    ADAAT-FGM    & 0.9443 & 0.9448 & 0.9444 & 0.9442 \\
    RRD-N$^{\dagger}$ & 0.9093 & 0.9094 & 0.9093 & 0.9092 \\
    MST           & 0.9325 & 0.9345 & 0.9327 & 0.9324 \\
    \midrule
    \method       & \textbf{0.9689} & \textbf{0.9698} &
                    \textbf{0.9691} & \textbf{0.9689} \\
    \bottomrule
  \end{tabular}
\end{table}

\begin{figure*}[t]
  \centering
  \includegraphics[width=0.96\textwidth]{figures/fig3_training_dynamics.pdf}
  \Description{Two line charts show validation Macro-F1 across epochs for the
  five PHEME folds and the five Ma-Weibo folds. Hollow circles mark the epochs
  from which the final checkpoints were selected.}
  \caption{Validation Macro-F1 trajectories of \method\ on the five outer
  folds. Hollow circles denote the selected checkpoint epochs; selection uses
  the validation-only multi-metric rule rather than Macro-F1 alone.}
  \label{fig:training-dynamics}
\end{figure*}

\paragraph{Ma-Weibo comparison.}
Table~\ref{tab:maweibo-results} shows that \method\ also performs strongly on
the larger Chinese corpus.  It improves accuracy from 0.9569 for the strongest
completed five-fold baseline, PAGCN, to 0.9689, an absolute gain of 1.20
percentage points.  The requested two-class aggregates show the same pattern:
macro-precision rises from 0.9570 to 0.9698, macro-recall from 0.9570 to
0.9691, and macro-F1 from 0.9569 to 0.9689.  Thus the gain is not an artifact
of reporting a support-weighted score on this nearly balanced corpus.
ADAAT-FGM and MST form the next strongest group, while the feature-engineered
and sequential baselines remain substantially lower.  The three completed
RRD-N folds average 0.9093 accuracy; because folds 4--5 are unavailable, we
retain it as an explicitly incomplete result and exclude it from five-fold
rankings and improvement calculations.

The absolute performance difference between the two datasets should not be
read as a language effect in isolation.  Ma-Weibo is nearly class-balanced and
contains substantially denser conversations, whereas PHEME contains 10.54\%
source-only events and a minority rumor class.  Accordingly, \method's margin
over the strongest completed baseline is larger on PHEME (2.97 points) than on
Ma-Weibo (1.20 points).  This pattern supports the intended role of unified
evidence learning under sparse and heterogeneous conversations, while avoiding
a causal claim from only two corpora.

\paragraph{Training dynamics.}
Figure~\ref{fig:training-dynamics} analyzes the optimization behavior of the
full model using the histories produced by the formal five-fold runs.
Validation Macro-F1 rises rapidly during the first four to six epochs and then
enters a narrower plateau on both datasets.  The selected epochs lie between
9 and 14 on PHEME and between 10 and 17 on Ma-Weibo, well before the 30-epoch
cap.  Continued optimization after the plateau produces small fold-specific
oscillations rather than a consistent upward trend.  This explains why a
moderate patience window is preferable to fixing the checkpoint at the maximum
epoch.

The checkpoint markers do not always coincide with the highest point on the
displayed Macro-F1 curve because model selection balances Accuracy,
Weighted-F1, Macro-F1, and Rumor-F1 on the validation split.  A checkpoint is
therefore retained only when its weakest validation measure is competitive,
rather than when one aggregate score peaks.  The similar rise-and-plateau
pattern across all ten folds shows that the selected models emerge from a
stable optimization regime rather than an isolated single-epoch spike.

The selected validation Macro-F1 values range from 0.8764 to 0.9147 on PHEME
and from 0.9679 to 0.9866 on Ma-Weibo.  More importantly, no fold exhibits a
late collapse after reaching its high-performance region: subsequent epochs
remain close to the plateau even when they are not selected.  The fold-to-fold
difference is therefore expressed mainly through the attainable plateau and
the epoch at which it is reached.  Thus the exact-replay SAM training runs
reach stable validation regions and support validation-based checkpointing
without requiring checkpoint averaging or post-hoc ensembling.

\begin{table}[t]
  \caption{Five-fold ablation on PHEME. $\Delta$Acc. is the absolute accuracy
  change in percentage points relative to the full model.}
  \label{tab:ablation}
  \centering
  \small
  \setlength{\tabcolsep}{3.1pt}
  \begin{tabular}{@{}lccccc@{}}
    \toprule
    Variant & Acc. & W-F1 & M-F1 & R-F1 & $\Delta$Acc. \\
    \midrule
    \method & \textbf{0.8988} & \textbf{0.8991} & \textbf{0.8925} &
      \textbf{0.8666} & -- \\
    w/o graph evidence & 0.8940 & 0.8942 & 0.8873 & 0.8601 & $-0.48$ \\
    w/o text evidence & 0.8601 & 0.8609 & 0.8525 & 0.8193 & $-3.88$ \\
    w/o cross-event context & 0.8985 & 0.8988 & 0.8923 & 0.8663 & $-0.03$ \\
    w/o view consistency & 0.8895 & 0.8902 & 0.8837 & 0.8579 & $-0.93$ \\
    w/o SAM & 0.8951 & 0.8955 & 0.8890 & 0.8629 & $-0.37$ \\
    \bottomrule
  \end{tabular}
\end{table}

\begin{figure}[t]
  \centering
  \includegraphics[width=\columnwidth]{figures/fig2_ablation_drops.pdf}
  \Description{A horizontal grouped bar chart showing the accuracy, macro-F1,
  and rumor-F1 decreases caused by five ablations. Removing text evidence has
  the largest decrease, followed by cross-view consistency, graph evidence,
  SAM, and cross-event context.}
  \caption{Absolute PHEME performance decrease after removing each component.
  Text evidence and cross-view consistency produce the largest drops; the
  cross-event context has little effect on aggregate five-fold metrics.}
  \label{fig:ablation-drops}
\end{figure}

\paragraph{Ablation study.}
Table~\ref{tab:ablation} and Figure~\ref{fig:ablation-drops} isolate the main
components under the same PHEME partitions and training seed.  Textual evidence
is indispensable: removing it lowers accuracy by 3.88 points and rumor-F1 by
4.73 points.  Cross-view consistency is the next largest contributor, with
0.93-point accuracy and 0.87-point rumor-F1 reductions, showing that agreement
between label-preserving serializations matters beyond merely sharing an
encoder.  Removing graph evidence reduces accuracy by 0.48 points and
rumor-F1 by 0.65 points, confirming a complementary but smaller structural
contribution.  Exact-replay SAM contributes 0.37 accuracy points without
changing inference.

The cross-event context ablation changes all four aggregate metrics by only
about 0.03 points.  Thus the present evidence does not support attributing the
overall gain primarily to the reference bank.  It remains a well-defined source
of support for individual queries, but establishing a selective benefit for
sparse or difficult events requires a dedicated subgroup study.  Reporting
this small aggregate effect is important: the principal measured improvements
come from jointly learned textual and structural evidence, cross-view
regularization, and sharpness-aware optimization.

The ablations should be interpreted as component removals rather than an
additive decomposition.  In particular, the 0.93-point consistency drop is
measured while the shared text encoder remains active, so it isolates the
benefit of constraining two valid conversation serializations rather than the
benefit of language modeling itself.  Likewise, the graph ablation retains
text and semantic context, explaining why its 0.48-point drop is smaller than
the gap between graph-only baselines and \method.  The SAM result has a
different operational implication: its 0.37-point gain costs a second
sharpness pass during optimization, but introduces neither new parameters nor
additional computation after training.  These distinctions prevent equating
the largest numerical ablation with the only useful module and clarify which
components affect representation, robustness, and deployment cost.

The class-sensitive ablation changes further distinguish the evidence paths.
Removing text reduces rumor-F1 by 4.73 points, exceeding its 3.88-point
accuracy loss; removing graph evidence similarly reduces rumor-F1 by 0.65
points versus a 0.48-point accuracy loss.  These two channels therefore make a
disproportionately visible contribution to the minority rumor class.  By
contrast, removing cross-view consistency lowers accuracy and rumor-F1 by
similar amounts (0.93 and 0.87 points), indicating a broader stabilization of
the fused decision rather than an exclusively class-specific effect.  The SAM
ablation is also nearly uniform across aggregate metrics.  This metric-level
pattern complements the magnitude ranking in Figure~\ref{fig:ablation-drops}:
text is the primary evidence source, graph structure supplies smaller
minority-sensitive support, and consistency and SAM mainly improve the
reliability of their joint optimization.

\FloatBarrier
\section{Conclusion}
\label{sec:conclusion}

We presented \method, a single-checkpoint framework that unifies propagation,
multi-view event text, and cross-event semantic support for social rumor
detection.  The model couples complementary graph readouts with a shared
DeBERTa-v3 encoder, aligns label-preserving conversation serializations, and
uses branch supervision to prevent evidence collapse.  Exact-replay SAM makes
the two sharpness evaluations operate on the same accumulated stochastic
batch, adding no inference-time computation.

Across fixed five-fold partitions, \method\ achieves 0.8988 accuracy on PHEME
and 0.9689 on Ma-Weibo, outperforming the strongest completed five-fold
baseline by 2.97 and 1.20 percentage points, respectively.  The PHEME ablation
identifies textual evidence and cross-view consistency as the dominant
contributors, with smaller complementary gains from propagation evidence and
SAM.  Cross-event context yields only a marginal aggregate improvement, so its
value should be tested on pre-defined sparse-thread and difficult-event groups
rather than inferred from the overall mean.  Future work should also reduce
the training cost of large language encoders and exact-replay SAM, and examine
generalization across platforms, topics, and evolving rumor distributions.

\section{Ethical Considerations}

Rumor detection is a high-impact classification task.  False positives may
suppress legitimate reporting, satire, or unresolved claims, while false
negatives may allow harmful misinformation to spread.  The model should
therefore support human review rather than autonomous censorship; its output
is a probabilistic signal based on textual and conversational patterns, not a
fact-checking verdict or proof of author intent.

The experiments use existing research datasets of public social-media
conversations and introduce no external user profiles or web evidence.
Released artifacts should replace user identifiers with event-local indices
and follow the original dataset terms.  Aggregate metrics may conceal unequal
error rates across topics, languages, communities, and conversation sizes, so
we report class-specific results and retain source-only events rather than
silently filtering them.  Targeted subgroup evaluation across topics,
communities, and sparse-thread strata remains necessary before deployment.

\bibliographystyle{../acmart/ACM-Reference-Format}
\bibliography{references}

\end{document}
